% ==================================================================
% id:       lim-012
% topic:    גבולות ורציפות של פונקציות
% subtopic: רציפות ונקודות אי־רציפות
% solution: full
% ==================================================================
\textbf{
בכל אחד מהסעיפים הבאים,
מצאו את  ערכי
$a,b$
עבורם הפונקציה רציפה בתחום הגדרתה.}
\begin{enumerate}

\item
$$\ds{f(x)=\case{ccc}{\lr{\frac{1+\tan x}{1+\sin x}}^{\frac{\ln 9}{\sin^3x}} & x<0 \\ ax+b & 0\leq x\leq 3 \\ \frac{2x^3-9x^2+27}{2x^3-17x^2+48x-45} & x>3}}$$
\hint{בדקו רציפות בנקודות התפר. ב-$0$ יש גבול מהצורה $1^\infty$, וב-$3$ המונה והמכנה מתאפסים ולכן יש לפרק אותם לגורמים.}
\begin{solbox}
נציג
\begin{align*}
    \lr{\frac{1+\tan x}{1+\sin x}}^{\frac{\ln 9}{\sin^3x}}&=\ \lr{\frac{1{\color{red}+\sin x-\sin x}+\tan x}{1+\sin x}}^{\frac{\ln 9}{\sin^3x}}\\
    &=\ \lr{1+\frac{\tan x-\sin x}{1+\sin x}}^{\frac{\ln 9}{\sin^3x}}.
\end{align*}
נשים לב
(מאריתמטיקה ורציפות)
כי
$$.\lim_{x\to0}\frac{\tan x-\sin x}{1+\sin x}=0$$
נמשיך לפשט
    $$,\lr{1+\frac{\tan x-\sin x}{1+\sin x}}^{\frac{\ln 9}{\sin^3x}}
    =\left[\lr{1+\frac{\tan x-\sin x}{1+\sin x}}^{\color{red}\frac{1+\sin x}{\tan x-\sin x}}\right]^{{\color{red}\cdot\frac{\tan x-\sin x}{1+\sin x}\cdot}\frac{\ln 9}{\sin^3x}}$$
הביטוי בבסיס מתכנס ל-
$e$,
לכן נבדוק מה גבולה של הפונקציה במעריך.
נציג
\begin{align*}
    \frac{\tan x-\sin x}{1+\sin x}\cdot\frac{\ln 9}{\sin^3x}&=\ \frac{\frac{\tan x}{\color{red}\sin x}-\frac{\sin x}{\color{red}\sin x}}{{1}+{\sin x}}\cdot\frac{\ln 9}{\frac{\sin^3x}{\color{red}\sin x}}\\
    &=\ \frac{\frac{1}{\cos x}-1}{1+\sin x}\cdot\frac{\ln 9}{\sin^2x}\\
    &=\ \ln 9\cdot\frac{1-\cos x}{\sin^2 x\cos x(1+\sin x)}\\
    &=\ \ln 9\cdot\frac{1-\cos x}{\color{red}x^2}\cdot\frac{\color{red}x^2}{{\sin^2 x}}\cdot\frac{1}{\cos x(1+\sin x)},
\end{align*}
מכאן
\begin{align*}
\lim_{x\to 0}\frac{\tan x-\sin x}{1+\sin x}\cdot\frac{\ln 9}{\sin^3x}&=\ \ln 9\cdot\lim_{x\to0}\frac{1-\cos x}{x^2}\cdot\left(\frac{x}{{\sin x}}\right)^2\cdot\frac{1}{\cos x(1+\sin x)}\\
&\overset{(\ast)}{=}\ \ln 9\cdot\frac{1}{2}\cdot1\cdot1\\
&=\ \ln3,
\end{align*}
לכן
$$.\lim_{x\to0^-}f(x)=\lim_{x\to0^-}\lr{\frac{1+\tan x}{1+\sin x}}^{\frac{\ln 9}{\sin^3x}}\overset{(\ast)}{=}\e^{\ln3}=3$$
מצד שני
$$.\lim_{x\to0^+}f(x)=\lim_{x\to0^+}(ax+b)=b$$
לכן לרציפותה של
$f$
ב-
$x=0$,
נדרוש כי
$b=3$.

$(\ast)$
השתמשנו כאן בגבולות המופלאים
$$.\lim_{t\to 0}\frac{\sin t}{t}=1,\ \lim_{t\to 0}\frac{1-\cos t}{t^2}=\frac{1}{2},\ \lim_{t\to 0}(1+t)^{\frac{1}{t}}=e$$

נעבור לדיון ברציפות של
$f$
בנקודה
$x=3$.
מתקיים
$$.\lim_{x\to3^-}f(x)=\lim_{x\to3^-}(ax+b)=3a+3$$
מצד שני נציג
\begin{align*}
    \frac{2x^3-9x^2+27}{2x^3-17x^2+48x-45}&\overset{(\ast)}{=}\ \frac{(x-3)^2(2x+3)}{(x-3)^2(2x-5)}\\
    &=\ \frac{2x+3}{2x-5},
    \end{align*}
לכן
\begin{align*}
    \lim_{x\to3^+}f(x)&=\ \lim_{x\to3^+}\frac{2x^3-9x^2+27}{2x^3-17x^2+48x-45}\\
    &=\ \lim_{x\to3^+}\frac{2x+3}{2x-5}\\
    &=\ 9.
\end{align*}
$(\ast)$
נשים לב כי הפולינומים במונה ובמכנה מתאפסים ב-
$x=3$,
לכן ניתן להציגם כ-
$(x-3)(ax^2+bx+c)$.
לדוגמא, עבור המונה
$$,2x^3-9x^2+27=(x-3)(ax^2+bx+c)$$
לאחר פתיחת סוגריים נקבל
$$.2x^3-9x^2+27=ax^3+(b-3a)x^2+(c-3b)x-3c$$
לאחר השוואת המקדמים נקבל כי
$$,a=2,\ -9=b-3a,\ c-3b=0,\ -3c=27$$
כלומר
$$.a=2,\ b=-3,\ c=-9$$
כך ננהג גם עם המכנה.

שיטה זו נקראת חילוק פולינומים, נעסוק בה בהרחבה בשלבים מתקדמים של הקורס.
לכן לרציפותה של
$f$
ב-
$x=3$
נדרוש ש-
$3a+3=9$,
כלומר
$a=2$.
לסיכום,
$$.a=2,\ b=3$$
\end{solbox}

\item $$\ds{f(x)=\case{ccc}{\frac{3x^2-9x-4}{x^2-1}-\frac{3x-5}{2x+2} & x<-1 \\ ax+b & -1\leq x\leq 0 \\ \lr{1+x\cdot\sin(2x)}^{\frac{4\ln 8}{\cos x-\cos (5x)}} & x>0}}$$
\hint{בדקו רציפות בנקודות התפר: משמאל ל-$-1$ אחדו את השברים וצמצמו, ומימין ל-$0$ יש גבול מהצורה $1^\infty$.}
\begin{solbox}
נציג
\begin{align*}
    \frac{3x^2-9x-4}{x^2-1}-\frac{3x-5}{2x+2}&=\ \frac{3x^2-10x-13}{2(x-1)(x+1)}\\
    &=\ \frac{(x+1)(3x-13)}{2(x-1)(x+1)}\\
    &=\ \frac{3x-13}{2(x-1)}.
\end{align*}
לכן
\begin{align*}
    \lim_{x\to-1^-}f(x)&=\ \lim_{x\to-1^-}\left(\frac{3x^2-9x-4}{x^2-1}-\frac{3x-5}{2x+2}\right)\\
    &=\ \lim_{x\to-1^-}\frac{3x-13}{2(x-1)}\\
    &=\ 4.
\end{align*}
מצד שני
$$.\lim_{x\to-1^+}f(x)=\lim_{x\to-1^+}(ax+b)=-a+b$$
לרציפותה של
$f$
בנקודה
$x=-1$
נדרוש
$-a+b=4$.

נעבור לדיון ברציפות בנקודה
$x=0$.
נציג
\begin{align*}
    \lr{1+x\cdot\sin(2x)}^{\frac{4\ln 8}{\cos x-\cos (5x)}}&=\ \lr{1+x\cdot\sin(2x)}^{\frac{4\ln 8}{\cos x-\cos (5x)}{\color{red}\cdot\frac{x\cdot\sin(2x)}{x\cdot\sin(2x)}}}\\
    &=\ \lr{1+x\cdot\sin(2x)}^{\frac{1}{{x\cdot\sin(2x)}}\cdot\frac{4\ln 8\cdot x\cdot\sin(2x)}{\cos x-\cos (5x)}}.
\end{align*}
נשים לב כי
\begin{align*}
    \frac{4\ln 8\cdot x\cdot\sin(2x)}{\cos x-\cos (5x)}&=\ 4\ln 8\cdot\frac{ x\cdot\sin(2x)}{\cos x{\color{red}-1+1}-\cos (5x)}\\
    &=\ 4\ln 8\cdot\frac{\frac{x}{\color{red}x}\cdot{\color{red}2\cdot}\frac{\sin(2x)}{\color{red}2x}}{\frac{\cos x-1}{\color{red}x^2}+{\color{red}25\cdot}\frac{1-\cos (5x)}{\color{red}(5x)^2}},
\end{align*}
לכן
\begin{align*}
    \lim_{x\to0}\frac{4\ln 8\cdot x\cdot\sin(2x)}{\cos x-\cos (5x)}&=\ 8\ln 8\cdot\lim_{x\to0}\frac{\frac{\sin(2x)}{2x}}{\frac{\cos x-1}{x^2}+{25\cdot}\frac{1-\cos (5x)}{(5x)^2}}\\
    &\overset{(\ast)}{=}\ \frac{2\ln 8}{3},
\end{align*}
מכאן
\begin{align*}
    \lim_{x\to0^+}f(x)&=\ \lim_{x\to0^+}\lr{1+x\cdot\sin(2x)}^{\frac{4\ln 8}{\cos x-\cos (5x)}}\\
    &=\ \lim_{x\to0^+}\left[\lr{1+x\cdot\sin(2x)}^{\frac{1}{x\cdot\sin(2x)}}\right]^{\frac{4\ln 8\cdot x\cdot\sin(2x)}{\cos x-\cos (5x)}}\\
    &\overset{(\ast)}{=}\ e^\frac{2\ln 8}{3}\\
    &=\ 4,
\end{align*}
מצד שני
$$.\lim_{x\to0^-}f(x)=\lim_{x\to0^-}(ax+b)=b$$
לכן לרציפותה של
$f$
ב-
$x=0$
נדרוש
$b=4$.
בסה"כ קיבלנו
$-a+b=4$
ו-
$b=4$,
כלומר
$$. a=0,\ b=4$$

$(\ast)$ יש כאן שוב שימוש בגבולות מופלאים
אשר צוטטו בסעיף הקודם.
הערה: בחישוב הגבול השתמשנו באריתמטיקה של גבולות: הגבול של סכום, של מכפלה ושל מנה של פונקציות מתכנסות שווה לסכום, למכפלה ולמנה של הגבולות, כל עוד הגבול של המכנה שונה מ-$0$.
\end{solbox}

\item $$\ds{f(x)=\case{cc}{\frac{x^2+ax+b}{x-1} & x>1 \\ e^{2x}-x^2 & x\leq 1}}$$
\hint{כדי שהגבול מימין ל-$1$ יהיה סופי, המונה חייב להתאפס ב-$1$. היעזרו בזה כדי לפרק את המונה לגורמים.}
\begin{solbox}
הפונקציה
$g(x)=\frac{x^2+ax+b}{x-1}$
רציפה לכל
$x\neq 1$,
הפונקציה
$h(x)=e^{2x}-x^2$
רציפה לכל
$x\in\mathbb{R}$.
לכן לרציפותה של
$f(x)$
יש לדרוש רציפות ב"נקודת תפר", דהיינו ב-
$x=1$.
נדרוש
$$.\lim_{x\to1^+}\frac{x^2+ax+b}{x-1}=\lim_{x\to1^+}f(x)=\lim_{x\to1^-}f(x)=\lim_{x\to1^-}e^{2x}-x^2$$
מתקיים
\begin{equation}\label{first-condition-ex7}
    \lim_{x\to1^-}(e^{2x}-x^2)=e^2-1
\end{equation}
נשים לב כי
$$,\lim_{x\to1^+}(x^2+ax+b)=1+a+b$$
אם
$1+a+b\neq 0$,
אז
$$\lim_{x\to1^+}\frac{x^2+ax+b}{x-1}=\infty$$
או
$$,\lim_{x\to1^+}\frac{x^2+ax+b}{x-1}=-\infty$$
בפרט,
הפונקציה הנתונה לא רציפה בנקודה
$x=1$.
לכן נדרוש כי
\begin{equation}\label{third-condition-ex7}
    a+b+1=0
\end{equation}
לכן
ניתן להציג
$$,x^2+ax+b=(x-1)(x-x_0)$$
לאחר השוואת המקדמים בשני האגפים נקבל כי
$$.x^2+ax+b=(x-1)(x-b)$$
מכאן נסיק כי
$$\frac{x^2+ax+b}{x-1}=\frac{(x-1)(x-b)}{x-1}=x-b$$
ולכן
\begin{equation}\label{second-condition-ex7}
    .\lim_{x\to1^+}\frac{x^2+ax+b}{x-1}=\lim_{x\to1^+}(x-b)=1-b
\end{equation}
לאחר איסוף הממצאים שב-
)\ref{first-condition-ex7}(, )\ref{third-condition-ex7}(, )\ref{second-condition-ex7}(,
נקבל מערכת משוואות
\[
,\begin{cases}
a+b=-1,\\
b=2-e^{2}.
\end{cases}
\]
אשר פתרונה הוא
$$.a=-3+e^2,\ b=2-e^2$$
הערה: בחישוב הגבול השתמשנו באריתמטיקה של גבולות: הגבול של סכום, של מכפלה ושל מנה של פונקציות מתכנסות שווה לסכום, למכפלה ולמנה של הגבולות, כל עוד הגבול של המכנה שונה מ-$0$.
\end{solbox}

\item $$\ds{f(x)=\case{cc}{\frac{\tan\lr{x^3}\cos\lr{\frac\pi3e^x}}{e^{x^3}-1} & x<0 \\ ax+b & 0\leq x\leq 2 \\ \frac{8x}{x^2-4}-\frac{8x+4}{2x^2-3x-2} & x>2}}$$
\hint{בדקו את נקודות התפר: ב-$0$ השתמשו בגבולות המופלאים של $\tan t$ ושל $e^t-1$, וב-$2$ אחדו את השברים וצמצמו.}
\begin{solbox}
הפונקציה
$g(x)=\frac{\tan\lr{x^3}\cos\lr{\frac\pi3e^x}}{e^{x^3}-1}$
רציפה לכל
$x\neq 0$,
הפונקציה
$h(x)=ax+b$
רציפה לכל
$x\in\mathbb{R}$,
והפונקציה
$k(x)=\frac{8x}{x^2-4}-\frac{8x+4}{2x^2-3x-2}$
רציפה לכל
$x\neq 2$
ו-
$x\neq-\frac{1}{2}$.
היות ו-
$f(x)=k(x)$
לכל
$x>2$,
אז הנקודה
$x=-\frac{1}{2}$
לא רלוונטית לדיון.
אנו רואים שדי לדרוש את הרציפות של
$f$
בנקודות תפר, זאת אומרת בנקודות
$x=0$
ו-
$x=2$.
נציג
    $$,\frac{\tan\lr{x^3}\cos\lr{\frac\pi3e^x}}{e^{x^3}-1}=\frac{\color{red}x^3}{e^{x^3}-1}\cdot\frac{\tan(x^3)}{\color{red}x^3}\cdot\cos(\tfrac{\pi}{3}e^x)$$
    לכן
\begin{align*}
    \lim_{x\to0^-}f(x)&=\ \lim_{x\to0^-}\frac{\tan\lr{x^3}\cos\lr{\frac\pi3e^x}}{e^{x^3}-1}\\&=\ \lim_{x\to0^-}\frac{x^3}{e^{x^3}-1}\cdot\frac{\sin(x^3)}{x^3}\cdot\cos(\tfrac{\pi}{3}e^x)\cdot\frac{1}{\cos(x^3)}\\
    &=\ 1\cdot1\cdot\frac{1}{2}\cdot1\\
    &=\ \frac{1}{2},
\end{align*}
מצד שני
$$.\lim_{x\to0^+}f(x)=\lim_{x\to0^+}(ax+b)=b$$
אז לרציפותה של
$f$
ב-
$x=0$
נדרוש
\begin{equation}\label{f-cond-ex7}
    .b=\frac{1}{2}
\end{equation}
נעבור לדיון ברציפות ב-
$x=2$.
נציג
\begin{align*}
    \frac{8x}{x^2-4}-\frac{8x+4}{2x^2-3x-2}&=\ \frac{8x}{(x-2)(x+2)}-\frac{8x+4}{2(x-2)(x+\tfrac{1}{2})}\\
    &=\ \cdots\\
    &=\ \frac{4}{x+2},
\end{align*}
לכן

    $$,\lim_{x\to2^+}f(x)=\lim_{x\to2^+}\frac{4}{x+2}=1$$
מצד שני
$$,\lim_{x\to2^-}f(x)=\lim_{x\to2^-}(ax+b)=2a+b$$
אז לרציפותה של
$f$
ב-
$x=2$
נדרוש
\begin{equation}\label{s-cond-ex7}
    .1=2a+b
\end{equation}
נאחד את התנאים
(\ref{f-cond-ex7})
ו-
(\ref{s-cond-ex7}),
ונקבל
$$,b=\frac{1}{2},\ 2a+b=1$$
כלומר
$$.a=\frac{1}{4},\ b=\frac{1}{2}$$
הערה: בחישוב הגבול השתמשנו באריתמטיקה של גבולות: הגבול של סכום, של מכפלה ושל מנה של פונקציות מתכנסות שווה לסכום, למכפלה ולמנה של הגבולות, כל עוד הגבול של המכנה שונה מ-$0$.
\end{solbox}

\item $$\ds{f(x)=\case{cc}{e^{ax} & x\geq 1 \\ \lr{2-x^2}^{\frac2{x-1}} & x\in[-1,1) \\ \ln\lr{-3x+b} & x<-1}}$$
\hint{בדקו את נקודות התפר: ב-$1$ יש גבול מהצורה $1^\infty$. שימו לב שגם תחום ההגדרה של הלוגריתם מגביל את $b$.}
\begin{solbox}
הפונקציה
$g(x)=e^{ax}$
רציפה
לכל
$x\in\mathbb{R}$,
הפונקציה
$h(x)=(2-x^2)^{\frac{2}{x-1}}$
רציפה בתחום
$[-1,1)$,
הפונקציה
$k(x)=\ln(-3x+b)$
רציפה
ב-
$(-\infty,-1)$
אם
$b>-3$.
בדומה לסעיפים הקודמים נבדוק את הרציפות בנקודות תפר:
$x=1$
ו-
$x=-1$.
מתקיים
$$,\lim_{x\to1^+}f(x)=\lim_{x\to1^+}e^{ax}=e^a$$
מצד שני
    $$\lim_{x\to1^-}f(x)=\lim_{x\to1^-}\lr{2-x^2}^{\frac2{x-1}}$$
נציג
\begin{align*}
    \lr{2-x^2}^{\frac2{x-1}}&=\ ({\color{red}1}-(x^2{\color{red}-1}))^{\frac{2}{x-1}}\\
    &=\ (1-(x^2-1))^{{\color{red}\frac{-1}{x^2-1}}\cdot\frac{2{\color{red}(1-x^2)}}{x-1}}\\
    &=\ (1-(x^2-1))^{{\frac{1}{x^2-1}}\cdot(-2)(x+1)},
\end{align*}
לכן
\begin{align*}
    \lim_{x\to1^-}f(x)&=\ \lim_{x\to1^-}(1-(x^2-1))^{{\frac{1}{x^2-1}}\cdot(-2)(x+1)}\\
    &=\ e^{-4}.
\end{align*}
מכאן נסיק כי
$a=-4$.
נעבור לדיון ברציפות בנקודה
$x=-1$.
מתקיים
$$,\lim_{x\to-1^+}f(x)=\lim_{x\to-1^+}(2-x^2)^{\frac{2}{x-1}}=1$$
מצד שני
$$,\lim_{x\to-1^-}f(x)=\lim_{x\to-1^-}\ln(-3x+b)=\ln(3+b)$$
נדרוש
$$,1=\ln(3+b),\ b> -3$$
כלומר
$$.b=e-3$$
לסיכום,
$$.a=-4,\ b=e-3$$
\end{solbox}

\item $$\ds{f(x)=\case{cc}{\frac{\sqrt[5]{1+2x+x^3}-\sqrt[5]{1+2x^2+x^3}}{ax-bx^2} & x\geq 0  \\ \frac{\sin^2(3x)}{1-\cos(x)} & x<0}}$$
\hint{הוציאו את אחד השורשים כגורם משותף והשתמשו בגבול המופלא $\lim_{t\to0}\frac{(1+t)^{\alpha}-1}{t}=\alpha$; משמאל ל-$0$ השתמשו בגבולות המופלאים של $\sin$ ו-$\cos$.}
\begin{solbox}
נתמקד כרגע בחקירת רציפות בנקודה
$x=0$.
נסמן
$$.g(x)=\frac{\sqrt[5]{1+2x+x^3}-\sqrt[5]{1+2x^2+x^3}}{ax-bx^2},\ h(x)=\frac{\sin^2(3x)}{1-\cos(x)}$$
נציג
\begin{align*}
    g(x)&=\ \sqrt[5]{1+2x^2+x^3}\cdot\frac{\frac{\sqrt[5]{1+2x+x^3}}{\sqrt[5]{1+2x^2+x^3}}-1}{ax-bx^2}\\
    &=\ \sqrt[5]{1+2x^2+x^3}\cdot\frac{\sqrt[5]{\frac{1+2x+x^3}{{1+2x^2+x^3}}}-1}{ax-bx^2}\\
    &=\ \sqrt[5]{1+2x^2+x^3}\cdot\frac{\left(\frac{1+2x+x^3}{{1+2x^2+x^3}}\right)^{\frac{1}{5}}-1}{ax-bx^2}\\
    &=\ \sqrt[5]{1+2x^2+x^3}\cdot\frac{\left(\frac{1{\color{red}+2x^2}+x^3+2x{\color{red}-2x^2}}{{1+2x^2+x^3}}\right)^{\frac{1}{5}}-1}{ax-bx^2}\\
    &=\ \sqrt[5]{1+2x^2+x^3}\cdot\frac{\left(1+\frac{2x-2x^2}{1+2x^2+x^3}\right)^{\frac{1}{5}}-1}{ax-bx^2}\\
    &=\ \sqrt[5]{1+2x^2+x^3}\cdot\frac{\left(1+\frac{2x-2x^2}{1+2x^2+x^3}\right)^{\frac{1}{5}}-1}{\color{red}\frac{2x-2x^2}{1+2x^2+x^3}}\cdot\frac{\color{red}\frac{2x-2x^2}{1+2x^2+x^3}}{ax-bx^2}\\
    &=\ \sqrt[5]{1+2x^2+x^3}\cdot\frac{\left(1+\frac{2x-2x^2}{1+2x^2+x^3}\right)^{\frac{1}{5}}-1}{\frac{2x-2x^2}{1+2x^2+x^3}}\cdot\frac{\frac{2-2x}{1+2x^2+x^3}}{a-bx},
\end{align*}
מכיוון ש-
$$\lim_{x\to0}\frac{2x-2x^2}{1+2x^2+x^3}=0$$
נסיק כי
$$,\lim_{x\to0}\frac{\left(1+\frac{2x-2x^2}{1+2x^2+x^3}\right)^{\frac{1}{5}}-1}{\frac{2x-2x^2}{1+2x^2+x^3}}=\frac{1}{5}$$
ולכן
\begin{align*}
    \lim_{x\to0^+}f(x)&=\ \lim_{x\to0}g(x)\\
    &=\ \lim_{x\to0}\sqrt[5]{1+2x^2+x^3}\cdot\frac{\left(1+\frac{2x-2x^2}{1+2x^2+x^3}\right)^{\frac{1}{5}}-1}{\frac{2x-2x^2}{1+2x^2+x^3}}\cdot\frac{\frac{2-2x}{1+2x^2+x^3}}{a-bx}\\
    &=\ 1\cdot\frac{1}{5}\cdot\frac{2}{a}\\
    &=\ \frac{2}{5a}.
\end{align*}
נטפל עתה בגבול משמאל.
נציג
    $$\frac{\sin^2(3x)}{1-\cos(x)}=\frac{\sin^2(3x)}{\color{red}(3x)^2}\cdot\frac{\color{red}x^2}{1-\cos x}{\color{red}\cdot9}$$
מכאן
\begin{align*}
\lim_{x\to0^-}f(x)&=\ 9\cdot\lim_{x\to0}\frac{\sin^2(3x)}{(3x)^2}\cdot\frac{x^2}{1-\cos x}\\
&=\ 9\cdot1\cdot2\\
&=\ 18.
\end{align*}
מכאן נדרוש כי
$$,18=\frac{2}{5a}$$
כלומר,
$$.a=\frac{1}{45},\ b\in\mathbb{R}$$
הערה: בחישוב הגבול השתמשנו באריתמטיקה של גבולות: הגבול של סכום, של מכפלה ושל מנה של פונקציות מתכנסות שווה לסכום, למכפלה ולמנה של הגבולות, כל עוד הגבול של המכנה שונה מ-$0$.
\end{solbox}
\end{enumerate}
